\chapter{Methodology and Experimental Design} \label{ch:methodology}
This chapter describes how the raw data was used to answer three questions, in the
order the Results chapter answers them. It defines the splits, feature sets, models,
targets and evaluation procedures. The same held-out observations were used whenever
models were compared.

\paragraph*{Experimental framework:} Q1 attributes a plot's persistent departure from
the shared growth curve to environmental and management conditions. Q2 tests whether
letting that relationship vary spatially explains what Q1's global models cannot. Q3,
supplementary, tests whether any of this signal also improves top-height prediction,
and reports what a physics-informed architecture did and did not achieve.
\begin{table}[!ht]
\centering
\footnotesize
\renewcommand{\arraystretch}{1.15}

\begin{tabular}{
@{}
>{\raggedright\arraybackslash}p{0.9cm}
>{\raggedright\arraybackslash}p{3.5cm}
>{\raggedright\arraybackslash}p{3.2cm}
>{\raggedright\arraybackslash}p{2.8cm}
@{}
}
\toprule
Q & Unit and target & Models or comparison & Main evaluation \\
\midrule
Q1 & Plot; mean pooled-CR residual & Elastic Net, XGBoost and two-stage
mixed-effects model & Five-fold spatial block CV; coefficients, SHAP and residual Moran's $I$ \\
Q2 & Plot; difference between fitted and yield-class-implied asymptotes & Elastic Net, XGBoost
and GNNWR & Five-fold spatial block CV; attribution, Moran's $I$ and target audit \\
Q3 & Plot--survey row; 95th-percentile top height & CR, linear regression, random forest,
XGBoost, DNN, PINN and PINN-$k$ & Five-fold spatial block CV; random-split and PINN
forward-pass diagnostics \\
Q3 (bridge) & Plot--survey row; error reduction from the fold-matched CR curve & Conditioned
DNN, PINN, PINN-$k$ and XGBoost; no-env and permutation controls & Five-fold spatial block
CV; CR-error quartiles \\
\bottomrule
\end{tabular}
\caption{Units, targets and evaluation used for each question. The Q3 bridge check reuses
Q3's own predictions to test whether environmental conditioning corrects the curve at row
level; it is reported as part of Q3, not as a fourth question.}
\label{tab:rq-framework}
\end{table}

\begin{figure}[!htbp]
\centering
\includegraphics[width=\textwidth]{figures/fig_methodology/method_fig_pipeline-overview.png}
\caption{\small{Analysis pipeline. The two cohorts feed shared spatial partitions and separate feature sets. Q1 examines departure from a shared curve, Q2 examines plot-specific asymptote deviation, and Q3 predicts top height.}}
\label{fig:pipeline-overview}
\end{figure}

\section{Data partitions and leakage control}
The four-survey cohort was the main dataset; the smaller six-survey cohort provided a
nested sensitivity analysis with longer observation histories. The unit of analysis
depended on the question. Q1 and Q2 aggregated the data to one observation per plot.
Q3 used plot--survey observations.

\subsection*{Evaluation splits}
Spatial-block cross-validation was the primary evaluation design. Whole forest compartments were assigned to folds, preventing the same compartment from appearing in both training and testing \cite{roberts2017Cross_C5_EVAL}. Training plots within 60\,m of a validation or test plot were also removed, while all held-out plots were retained. Five folds were balanced approximately by row count. In each rotation, one fold was used for testing, the following fold for validation and the remaining three for training. Each compartment was therefore tested once and predictions from the five test folds were pooled before calculating the overall metrics \cite{ploton2020Spatial_C5_EVAL}.

A random plot-level split assigned plots to training, validation and testing in a 60/20/20 ratio. This provided an easier reference without spatial separation, used in Section~\ref{sec:results-rq1-cr-reduction} to show how much a non-spatial split inflates apparent accuracy. Random assignments used seed~42. \limit{A temporal split (holding out later survey years) was also constructed for Q3 but is not reported: the chapter's scope was narrowed to the spatial-block comparison and the Q1/Q2 attribution questions, and the temporal holdout was not carried through to a full, comparably-robust analysis. Figure~\ref{fig:splits-combined} should be re-cropped to show only panels A and B before submission.}

\begin{figure}[!htbp]
\includegraphics[width=\textwidth]{figures/fig_methodology/method_fig-splits-combined_draft.png}
\caption{Evaluation data-splitting strategies: (A) random plot-level split and (B) compartment-level spatial-block split with a 60\,m buffer.}
\label{fig:splits-combined}
\end{figure}

Continuous predictors were standardised using parameters estimated from the training partition. The same transformation was applied to validation and test data. Age was scaled separately so that the PINN derivative could be converted to metres per year. Categorical soil classes were one-hot encoded. Compartment, sub-compartment and block identifiers were excluded from the model inputs. Two neighbour-height variables were also excluded because they had been calculated from nearby plots' observed heights before partitioning and could therefore leak held-out outcomes.

%------------------------------------------------
\section{Environmental feature sets}

Because Q1, Q2, and Q3 target different response variables, feature sets are constructed separately for each question using a shared four-step protocol:

\begin{enumerate}
\setlength\itemsep{0pt}
\item \textbf{Coverage screening:} Variables with insufficient spatial coverage are removed prior to ranking.
\item \textbf{Collinearity filtering:} For collinear pairs with Spearman $|\rho| \ge 0.95$, one representative is retained, prioritising direct external measurements over derived indices.
\item \textbf{Composite ranking:} Remaining candidates are sorted by their average rank across univariate Spearman $|\rho|$, permutation feature importance, and drop-column loss in held-out $R^2$.
\item \textbf{Multicollinearity pruning:} Features are pruned iteratively until all variance inflation factors (VIF) fall below 5 \cite{obrien2007Caution_C5_EVAL}. Protected stand variables remain fixed in the matrix throughout this process to constrain the VIF of competing predictors.
\end{enumerate}

Four nested feature sets are defined from these steps. Set~1 contains only the protected baseline variables: canopy cover, time since thinning, its missingness indicator, and recent thinning. Set~2 adds the ten highest-ranked candidates. Set~3 expands on Set~2 by adding the top half of terrain and wind features, while Set~4 adds the top half of climate, soil, and spatial context features to Set~3. Appendix~\ref{app:feature-sets} lists complete variable assignments and VIF diagnostics.

Feature selection is performed once on the designated training partition rather than re-run across individual folds, which ensures all models train on identical feature subsets, meaning fold-to-fold variance reflects true model stability and spatial partitioning rather than shifting predictor selection.

%------------------------------------------------
\section{Q1: Persistent curve-departure attribution}
\label{sec:rq2b-method}
\label{sec:q1-method}

Q1 focuses on the four-survey cohort, as the six-survey cohort contains too few compartments for compartment-level spatial attribution. For each plot-survey observation, the deviation from expected height is:
\begin{equation}
r_{it}=H_{it}-H(a_{it}),
\end{equation}
yielding the time-averaged plot target:
\begin{equation}
\bar r_i=\frac{1}{T_i}\sum_{t=1}^{T_i}r_{it}.
\label{eq:mean-cr-residual}
\end{equation}
Positive values identify stands that consistently outperform the pooled Chapman--Richards curve (the pooled fit itself is defined in Section~\ref{sec:rq1-method}). Stand age is excluded from the candidate predictor matrix because it already defines the baseline expectation.

Elastic Net and XGBoost are evaluated under five-fold compartment cross-validation. Elastic Net tunes its penalty parameter and $\ell_1$ mixing ratio through nested cross-validation on each training partition. XGBoost uses a fixed configuration -- tree depth 4, learning rate 0.04, up to 500 boosting rounds with early stopping -- without further retuning; this is the same fixed configuration Q2's XGBoost uses, distinct from Q3's per-target tuned configuration (Section~\ref{sec:rq1-method}). Model performance is measured using pooled out-of-fold $R^2$.

To account for hierarchical spatial grouping, a two-stage mixed-effects model is also estimated. Stage~1 computes $\bar r_i$ from the non-linear baseline curve. Stage~2 fits:
\begin{equation}
\bar r_i=\boldsymbol{\beta}^{\top}\mathbf{x}_i+u_{c[i]}+\varepsilon_i,
\label{eq:mixed-effects}
\end{equation}
where $u_{c[i]}\sim\mathcal{N}(0,\sigma_u^2)$ represents a random intercept for each compartment, without random slopes. Fixed-effect coefficients are estimated using restricted maximum likelihood (REML) within each training fold. Spatial variance explained is calculated as the proportional reduction in compartment random-intercept variance relative to an intercept-only baseline; this comparison switches to maximum likelihood (ML) rather than REML for both models, since REML variance-component estimates are not cleanly comparable across models with different fixed-effect sets -- REML profiles out the fixed effects in a way that depends on how many there are, so a naive REML-vs-REML comparison would confound "the fixed effects explained variance" with "the REML transform changed." Parameters are reported as the mean $\pm$ sample standard deviation across folds.

Four diagnostic checks aid interpretation. First, Set~4 is re-estimated without canopy cover. Second, Pearson correlations are calculated between canopy cover and each target variable. Third, SHAP values are recomputed without canopy cover to determine whether correlated predictors absorb its importance. Fourth, between-compartment variance shares are computed for slope and TOPEX to resolve differences between the mixed model and Elastic Net.

Together, Elastic Net coefficients, mixed-effects fixed effects, and XGBoost mean absolute SHAP values provide three distinct views of growth attribution. SHAP calculations are performed post-estimation and do not influence feature selection. Finally, Moran's $I$ is computed on the pooled held-out residuals from Elastic Net and XGBoost to check for remaining spatial autocorrelation (method in Section~\ref{sec:eval-robustness}).

\ai{Proposed additional evaluation, not yet run: every Q1 number above (coefficients, $R^2$, Moran's $I$) reflects a single spatial-block split-seed (42). Q2's GNNWR is already reseeded three times to test sign stability (Section~\ref{sec:q2-method}); repeating that reseed for Q1's Elastic Net and XGBoost, reusing the same five-fold compartment assignments under different seeds, would test whether the slope/TOPEX stability contrast reported in Section~\ref{sec:results-q1-attribution} holds up outside one particular fold assignment, directly answering a gap already flagged in this project's own audit notes rather than leaving it a single-seed number reported as if stable.}

%------------------------------------------------
\section{Q2: Spatial attribution}
\label{sec:rq3-method}
\label{sec:q2-method}

Q1's global attribution models leave real, significant spatial autocorrelation in their residuals across every set, though Elastic Net's own residual clustering barely responds to added environmental data while XGBoost's responds substantially more (Section~\ref{sec:results-q1-attribution}), and a compartment-level random intercept only partially reduces it. This motivates testing whether letting the environment--deviation relationship itself vary spatially, rather than fitting one global relationship and shifting only its intercept locally, captures structure the global models cannot. GNNWR is used for this rather than classical GWR because it removes the need to pre-specify a distance kernel and bandwidth, and rather than a geographically weighted random forest because it keeps one interpretable linear coefficient per predictor at each location, preserving direct comparability with the Elastic Net and mixed-effects coefficients reported for Q1 (Section~\ref{sec:background-spatial-eval}).

Q2 estimates a distinct asymptotic height ceiling for each plot while fixing the Forestry Commission rate and shape parameters. For plot $i$ at survey time $t$, let
\begin{equation}
g_{it}=\left(1-e^{-p_{4,i}a_{it}}\right)^{p_{5,i}}.
\end{equation}
Fitting the single-parameter growth curve through the origin by ordinary least squares yields
\begin{equation}
\hat y_{\max,i}=\frac{\sum_t H_{it}g_{it}}{\sum_t g_{it}^2}.
\label{eq:plot-ymax-fit}
\end{equation}
The expected asymptote implied by standard yield class is calculated directly from the regional Forestry Commission parameters:
\begin{equation}
y_{\max,i}^{\mathrm{yldc}}=\mathrm{yldc}_i p_{2,i}+p_{1,i}+2p_{3,i}.
\end{equation}
The final attribution target is the deviation between these two ceilings:
\begin{equation}
\Delta y_{\max,i}=\hat y_{\max,i}-y_{\max,i}^{\mathrm{yldc}}.
\label{eq:local-ymax-diff}
\end{equation}
This formulation keeps the target strictly free of environmental covariates. Fixing $p_4$ and $p_5$ avoids over-parameterising curves from only four to six time points, though it absorbs any local shape mismatch into the fitted asymptote -- the mechanism behind the extreme-asymptote pattern audited below. Note $p$ here plays a different role from the pooled Chapman--Richards curve's own shape parameter fitted in Section~\ref{sec:rq1-method}: this section's $p_4$/$p_5$ are fixed, externally-sourced Forestry Commission values, not re-fitted quantities, despite the shared notation.

Elastic Net and XGBoost models are trained on $\Delta y_{\max}$ using feature Sets~2--4, with the same training setup as Q1 (Section~\ref{sec:q1-method}). To model spatially varying associations, Geographically Neural Network-Weighted Regression (GNNWR) is also estimated \cite{du2020Geographically_C6_GWR}:
\begin{equation}
\hat y_i=\beta_0(s_i)+\sum_j\beta_j(s_i)x_{ij}.
\label{eq:gnnwr}
\end{equation}
For each target location $s_i$, the spatial weighting sub-network receives a MinMax-scaled vector of Euclidean distances to all training reference plots. The network predicts multiplicative scaling factors that adjust a frozen, global ordinary least squares baseline for the intercept and each slope parameter. All reported runs use the full, uncapped training reference set without predefined distance kernels or fixed bandwidths. Predictions for held-out validation plots reference these same training points, evaluated across the shared five spatial folds.

\subsection*{Implementation constraints and simplifications}
Two separate hardware constraints shaped how GNNWR was fitted, disclosed here because they affect what the reported diagnostics can and cannot show. First, the spatial-weighting sub-network's input width equals the reference-set size, since it takes each plot's full distance vector to every reference point; at this project's scale ($\sim$31,000 training plots) this exhausted the available cluster GPU's memory on the first optimiser step. An initial workaround capped the reference set to a fixed row count, sampled proportionally across compartments so a smaller reference set still covered the forest geographically rather than a lucky or unlucky random slice; this was later superseded by removing a second, unrelated bottleneck (below), which allowed the reference set to be restored to its full, uncapped size for every headline result reported in this chapter. A separate reference-density ablation (Section~\ref{sec:results-q2-spatial}) deliberately re-introduces this cap at several sizes to test whether reference sparsity explains the six-survey cohort's instability; those runs are the one exception to "full, uncapped" and are reported as such. Second, and independently, the underlying GNNWR package's diagnostics class builds a classical-GWR hat matrix to report training-time AIC and F-test statistics, an $O(n^2 \times \text{features})$ tensor that exceeded available GPU memory regardless of reference-set size. This hat-matrix construction was disabled above a row-count threshold. It does not affect the $R^2$/RMSE values reported throughout this chapter, which are computed by a separate, direct formula; it does mean AIC/AICc are approximated by plain feature count rather than the true GWR-corrected effective degrees of freedom, and that the F1/F2/F3 significance tests classical GWR software reports are unavailable here. Neither constraint affects EN or XGBoost, which do not share GNNWR's architecture.

Predictor importance is assessed using both tree split gain and post-hoc SHAP values. XGBoost gain sums the training loss reduction across all tree splits involving a feature, reflecting its utility during model construction. In contrast, mean absolute SHAP values quantify each feature's average additive contribution to final predictions \cite{lundberg2017unified}. SHAP values are computed out of fold across each test partition and pooled before interpretation. A separate check groups predictors into terrain, wind, soil, climate, and spatial-position/edge-effect categories and measures each category's contribution by grouped permutation importance (mean $R^2$ drop when every column in that category is shuffled together) and by the change in residual Moran's $I$ when that category is removed entirely and the model refitted, on a single spatial-block split of Set4 (4survey only).

\subsection*{Target and outlier audit}
To prevent flawed time series from corrupting the analysis, a plot is excluded from the curve-fitting population if any consecutive survey interval meets either of two filtering rules:
\begin{itemize}
\setlength\itemsep{2pt}
\item \textbf{Clearfell pattern:} Top height falls by at least 25\%, canopy cover drops by at least 70\%, and standing volume drops by at least 70\%.
\item \textbf{Measurement inconsistency:} Top height falls by at least 25\%, while canopy cover drops by no more than 10\% and standing volume does not decrease.
\end{itemize}
Intervals exhibiting large height drops that match neither criterion are classified as ambiguous and kept in the dataset.

Unusually low or high fitted asymptotes are screened separately using Tukey's $1.5\,\mathrm{IQR}$ fences \cite{tukey1977Exploratory}:
\begin{equation}
\left[Q_1 - 1.5\,\mathrm{IQR},\; Q_3 + 1.5\,\mathrm{IQR}\right],
\label{eq:tukey-fences}
\end{equation}
calculated independently within each cohort. Similar non-parametric thresholds have been used to identify parameter estimation errors in non-linear growth modelling \cite{markovic2014NewApproach}. Because this fence captures distributional extremes rather than demonstrably erroneous records, flagged plots remain in the primary analysis. A sensitivity check recalculates pooled held-out metrics after removing these observations from saved predictions, without refitting the models.

The audit evaluates whether specific plots recur within the worst 1\% and 5\% of residuals, independently across Elastic Net, XGBoost, and GNNWR. It compares their within-plot residual range, GNNWR local coefficients, edge distance to compartment boundaries, sub-compartment clustering, and topographic wind shelter against the wider sample. To test a known survey-boundary issue, the worst-fitting survey year is identified, and each flagged plot is refitted after excluding its 2008 observation. These diagnostics assess whether errors originate from target construction rather than genuine ecological factors.

%------------------------------------------------
\section{Q3: Prediction and physics guidance}
\label{sec:rq1-method}
\label{sec:q3-method}

Q3 is supplementary: it returns to whether the attribution signal in Q1 and Q2 also
improves top-height prediction, and what a physics-informed architecture does with it.
It compares a biological curve, standard tabular models and neural models. The headline
comparison uses Set~3 and the same five spatial folds. The other feature sets test whether adding
environmental inputs changes the result.

\subsection{Chapman--Richards (CR) baseline}
\begin{equation}
H(a) = y_{\max}\left(1 - e^{-ka}\right)^p,
\label{eq:cr}
\end{equation}
where $H(a)$ is top height at age $a$, $y_{\max}$ is the asymptotic height ceiling, $k$ is the growth rate parameter, and $p$ governs curve shape. All three parameters ($y_{\max}$, $k$, $p$) are fitted freely by non-linear least squares across five initial starting points, retaining the solution with the lowest sum of squared errors -- unlike Q2's plot-level fit, which fixes the equivalent rate and shape parameters at Forestry Commission values (Section~\ref{sec:q2-method}). Parameters are calibrated exclusively on the training partition of each split and fold, then frozen. This prevents validation and test observations from influencing the baseline fit or the PINN loss terms derived from it.

\subsection*{Conventional models}
Three classical baselines represent different levels of complexity: ordinary linear regression, a bagged random forest, and gradient-boosted decision trees via XGBoost \cite{chen2016XGBoost_C4_NN}. All conventional models use the same Q3 feature sets and cross-validation folds.

For Q3, XGBoost hyperparameters were tuned specifically on the top-height target. Evaluating 27 candidate configurations across five folds and two cohorts produced 270 fits. Model selection relied strictly on validation $R^2$, keeping test sets untouched until final evaluation. Appendix~\ref{app:hyperparameters} lists the full search grid and chosen values. This tuned configuration is specific to Q3; Q1 and Q2's attribution models use the fixed configuration given in Section~\ref{sec:q1-method} instead, without further retuning.

\subsection{DNN and physics-informed models}
All neural architectures share a three-layer prediction trunk with 128 units per layer and Leaky-ReLU activations. Training uses the Adam optimiser with a base learning rate of $10^{-4}$, weight decay of $10^{-5}$, gradient clipping at 1.0, batch size 256, and early stopping with a patience of 40 epochs (up to 500 epochs total). A ReduceLROnPlateau scheduler (decay factor 0.8, patience 15 epochs) lowers the learning rate when validation loss stalls, remaining active across most of the optimisation trajectory. The standard DNN minimises mean squared error:
\begin{equation}
\mathcal{L}_{\mathrm{data}}=\frac{1}{N}\sum_{i=1}^{N}
\left(\hat H_i-H_i\right)^2.
\label{eq:data-loss}
\end{equation}

The environment-conditioned DNN concatenates the environmental variables to the main input. The PINN variants use the same prediction trunk, but environmental variables enter small auxiliary networks that shape the physics and trajectory losses (Equations~\ref{eq:physics-loss}--\ref{eq:trajectory-loss}), not the trunk's prediction itself. The architectures therefore share the same core capacity, while the environmental information reaches training through a separate pathway from the DNN's.

\limit{This separation between the trunk's prediction and the auxiliary networks was not designed to isolate a growth-ceiling-specific effect; it reflects a limitation in the implemented forward pass, identified after training and discussed in Section~\ref{sec:rq1-pinn-limitation}.} One consequence is diagnostically useful regardless of intent: because terrain never reaches the trunk, there is nothing terrain-derived for a random split to leak between train and test, which separates this architectural property from the standard DNN baseline's split sensitivity (Section~\ref{sec:results-rq1-cr-reduction}).

The analytical derivative of the CR curve is:
\begin{equation}
H'(a)=y_{\max}pk e^{-ka}\left(1-e^{-ka}\right)^{p-1}.
\label{eq:cr-derivative}
\end{equation}
The physics loss penalises discrepancies between this theoretical growth rate and the network derivative computed via automatic differentiation:
\begin{equation}
\mathcal{L}_{\mathrm{phys}}=\frac{1}{N}\sum_{i=1}^{N}
\left(\frac{\partial\hat H_i}{\partial a_i}-H'(a_i)\right)^2.
\label{eq:physics-loss}
\end{equation}
To maintain consistent units ($\mathrm{m\,yr^{-1}}$), the network derivative is divided by the training-set age standard deviation before computing the loss.

The trajectory loss uses consecutive training surveys from the same plot:
\begin{equation}
\mathcal{L}_{\mathrm{traj}}=\frac{1}{M}\sum_{j=1}^{M}
\left(\frac{\hat H_{j,2}-\hat H_{j,1}}{\Delta a_j}-H'(\bar a_j)\right)^2,
\label{eq:trajectory-loss}
\end{equation}
where $\bar a_j$ is the midpoint age. Pairs are formed only when both rows belong to the training
partition. The total loss is
\begin{equation}
\mathcal{L}=\mathcal{L}_{\mathrm{data}}+
\lambda_{\mathrm{phys}}\mathcal{L}_{\mathrm{phys}}+
\lambda_{\mathrm{traj}}\mathcal{L}_{\mathrm{traj}}+
10^{-5}\sum_{\theta}|\theta|.
\label{eq:total-loss}
\end{equation}
The default uses $\lambda_{\mathrm{phys}}=\lambda_{\mathrm{traj}}=1$. Joint weights 0, 1 and 2
are compared in the ablation. Weight 0 is the physics-disabled control.

The first conditioned PINN changes the asymptote:
\begin{equation}
y_{\max}^{(i)}=y_{\max}+f_y(\mathbf{x}_i).
\label{eq:ymax-adjust}
\end{equation}
PINN-$k$ also changes the rate while keeping it positive:
\begin{equation}
k^{(i)}=k\exp\!\left(f_k(\mathbf{x}_i)\right).
\label{eq:k-adjust}
\end{equation}
Both auxiliary networks have one 16-unit hidden layer and start near zero. The shape parameter
$p$ remains fixed in every variant.

The trunk's prediction never evaluates Equation~\ref{eq:cr} with $y_{\max}^{(i)}$ or $k^{(i)}$;
these enter only the physics and trajectory losses (Equations~\ref{eq:physics-loss}--
\ref{eq:trajectory-loss}). Environmental information therefore reaches the trunk only indirectly,
through those loss gradients, never as a direct input to the prediction itself. \limit{This is a
limitation of the implemented forward pass rather than a chosen ablation design; see
Section~\ref{sec:rq1-pinn-limitation}.}

\subsection*{Does environmental conditioning correct the curve at row level?}
\label{sec:rq2a-method}
This is a row-level, supplementary check, not this chapter's headline attribution result: it asks whether environmental conditioning corrects a plot--survey's error at all, as a bridge into Q1 and Q2's plot-level attribution models above. It evaluates whether environmental predictors explain departures from standard yield curves using held-out predictions from Q3's own models. For each test observation $i$, the absolute error reduction relative to the fold-matched Chapman--Richards baseline is
\begin{equation}
d_i=\left|H_i-\hat H_i^{\mathrm{CR}}\right|-
\left|H_i-\hat H_i^{\mathrm{model}}\right|.
\label{eq:residual-reduction}
\end{equation}
A positive value indicates that the model improves on the baseline prediction. The analysis reports mean $d_i$, its percentage of mean Chapman--Richards error, and the percentage of observations where $d_i>0$.

Within each held-out fold, test observations are ranked by absolute baseline error and divided into quartiles (Q1 to Q4). Q1 contains observations where the Chapman--Richards curve fits best, while Q4 contains those with the largest departures. Mean error reduction and the percentage of improved plots are calculated for each quartile. This is repeated for the environment-conditioned DNN, PINN and PINN-$k$ across Sets~2--4 and both cohorts, and for XGBoost on Set~3 and both cohorts. The Set~3 DNN is evaluated across five training seeds to measure parameter initialisation spread. Other configurations use seed 42, meaning their reported spread reflects fold-to-fold variation rather than seed variance.

Two controls, both run on XGBoost with Set~3 only, isolate the source of any improvement. First, XGBoost is evaluated with and without environmental inputs. Second, the environmental columns are permuted using a single shared row permutation across every row of a held-out fold (train, validation, and test alike), and XGBoost is refitted from scratch on the permuted data. This tests whether the performance gain requires real environment-plot linkage, avoiding the confound of evaluating an already-trained model on out-of-distribution inputs. The same error-reduction calculation is then repeated on the refitted model's held-out predictions. For the DNN and XGBoost (both Set~3, both cohorts), the spatial pattern of $d_i$ is assessed with Moran's $I$ using the same weights as the residual analysis.

%------------------------------------------------
\section{Evaluation and robustness}
\label{sec:eval-robustness}

For residual $e_i=H_i-\hat H_i$, the main metrics are
\begin{align}
\mathrm{RMSE}&=\sqrt{\frac{1}{N}\sum_i e_i^2}, &
\mathrm{MAE}&=\frac{1}{N}\sum_i|e_i|, \\
R^2&=1-\frac{\sum_i e_i^2}{\sum_i(H_i-\bar H)^2}, &
\mathrm{Bias}&=\frac{1}{N}\sum_i e_i.
\end{align}
Positive bias denotes under-prediction. Each table states whether it reports a fold mean or a
pooled metric. Fold summaries are the mean $\pm$ sample SD across five folds. Pooled metrics are
calculated once from all out-of-fold predictions. These are not interchangeable.

For the four-survey spatial analyses, pooled $R^2$ 95\% confidence intervals use 1,000 cluster bootstrap resamples of whole compartments. The 2.5th and 97.5th percentiles form the interval. This preserves spatial clustering better than resampling individual rows. The smaller six-survey Q2 results are reported as point estimates; GNNWR is also repeated over three split seeds to test sign stability.

A semivariogram measures how quickly residuals become less similar as distance between plots increases; its range is the distance beyond which they are no longer spatially correlated. Each model's range is fitted separately from its own residuals, searched up to 5000m by default, widened to 10,000m for cells that did not resolve within 5000m; two models' Moran's $I$ are not necessarily tested at the same distance, and any report of a Moran's I value states the distance it was computed at for this reason.

Global Moran's $I$ is calculated from held-out residuals using distance-band weights, with the threshold set by this semivariogram range, and significance assessed via 999 permutations (ranges used are tabulated in the Appendix):

\begin{equation}
I=\frac{N}{W}
\frac{\sum_i\sum_j w_{ij}(e_i-\bar e)(e_j-\bar e)}
{\sum_i(e_i-\bar e)^2}.
\label{eq:morans-i}
\end{equation}
It is used after fitting, not for feature selection.
Three outcomes occur: the range resolves within the search window (Moran's $I$ reported directly); the semivariance stays flat throughout the window (no spatial structure detected, Moran's $I$ not computed); or the semivariance is still rising at the window's edge (Moran's $I$ computed at that capped distance, reported as provisional, not equivalent to a resolved result). The same rule is applied to Q3's bridge check, Q1's residuals and the Elastic Net, XGBoost and GNNWR residuals in Q2.

Neural training variation is assessed over five seeds for the selected Q3 models. XGBoost and the
DNN are also compared fold by fold on identical test folds. The physics-weight ablation tests
weights 0, 1 and 2. For PINN-$k$, parameter separability is assessed by the pooled correlation
between learned $y_{\max}$ and $k$; its confidence interval uses a separate 2,000-resample compartment
bootstrap (not the 1,000-resample procedure used for the $R^2$ intervals above). These checks distinguish predictive accuracy from parameter stability.

\section{Reproducibility}

Split assignments, feature lists and model settings are saved with each run. Seed 42 is used unless
a stated reseed is part of the analysis. Full feature membership, hyperparameter searches,
secondary split results and detailed outlier tables are retained in the Appendix. Methodological
limitations and their effect on interpretation are evaluated after the results, before the final
conclusions.

%=======================================================================================
%=======================================================================================
